\documentclass[conference]{IEEEtran}
\usepackage{amsmath,amssymb,amsthm}
\usepackage{graphicx}
\usepackage{booktabs}
\usepackage{cite}
\usepackage[hidelinks]{hyperref}
\graphicspath{{../results/}}

\newtheorem{proposition}{Proposition}
\newcommand{\Pperp}{\mathbf{P}_{\perp}}
\DeclareMathOperator{\tr}{tr}
\DeclareMathOperator{\diag}{diag}

\begin{document}

\title{Receiver-Informed Localization of Resource-Selective Jamming in the 5G NR Uplink}

\author{\IEEEauthorblockN{Erhan Karakoca}
\IEEEauthorblockA{Department of Computer Engineering, Bo\u{g}azi\c{c}i University, Istanbul, Turkey}}

\maketitle

\begin{abstract}
A jammer that knows the 5G NR numerology can concentrate its energy on selected resource
blocks and OFDM symbols, for example only on the demodulation reference signal (DMRS). We
study how the gNB can output, for every uplink slot, a map of the attacked PRB-by-symbol cells
from quantities its receiver already computes: received power, decision-directed error vector
magnitude (EVM), and a DMRS residual obtained by projecting least-squares pilot estimates
outside the delay subspace of physical channels. We show that the residual's distribution
without jamming depends on neither the channel nor the SNR, which yields an analytic
constant-false-alarm threshold and closed-form detection probabilities, and that EVM alone
reports pilot jamming at the wrong symbols because the damage appears in unjammed data. A
66k-parameter U-Net fusing the three features reaches a cell-level F1 of 0.95 at a 1\%
per-cell false-alarm rate in a Sionna PUSCH simulation, and a 7.7k-parameter version (F1 0.94)
runs in 0.43\,ms per slot on one CPU core as a frozen graph. Used by the receiver to weight pilots and log-likelihood
ratios by the estimated jamming power, the map restores the link under strong jamming without
the penalty that hard erasure pays at weak jamming.
\end{abstract}

\begin{IEEEkeywords}
5G NR, jamming detection, physical-layer security, DMRS, CFAR, U-Net, Sionna.
\end{IEEEkeywords}

% ---------------------------------------------------------------------------------------------
\section{Introduction}

Jamming remains one of the simplest attacks on cellular links~\cite{pirayesh2022jamming,
lichtman20185g}. In 5G NR the attacker does not need to jam the whole carrier: knowing the
numerology and the position of the DMRS, it can hit only the pilots, only the data symbols, or
a band of resource blocks~\cite{lichtman20185g,flores2023implementation}, and with access to
the downlink control information it can follow the user's scheduled
PRBs~\cite{wang2026downlink,alaimo2025exploiting}. Pilot jamming is particularly efficient,
because a corrupted channel estimate spoils every symbol of the
slot~\cite{clancy2011efficient}.

Most jamming detectors decide per slot or per block, from key performance indicators or
spectrum data~\cite{kilinc2022jade,neasamonisanthi2025interfo,reusmuns2024senseoran}, or from
the eigenstructure of multi-antenna pilot blocks~\cite{mehrabian2026robust}. Time--frequency
resolution has been obtained from EVM per symbol and resource block~\cite{ornek2023securing,
ornek2024advanced}, so far as images without a detection rule, and from per-tone power tests
in robust OFDM receivers~\cite{li2012robust}. A gNB, however, needs to know \emph{which}
resources were attacked: to report and react to the attack, and to decode through it. Two
difficulties appear at this resolution. First, the place where the jammer transmits and the
place where its effect shows up differ: under pilot jamming the EVM rises on the
\emph{unjammed} data symbols, through the corrupted channel estimate. Second, a cell-level
detector runs thousands of tests per slot, so its false-alarm rate must be controlled at every
SNR and channel.

This paper makes the following contributions.
\begin{itemize}
\item We formulate jamming localization in the NR uplink as a per-slot binary map over PRBs
  and OFDM symbols whose labels mark the \emph{attacked} cells, and show that EVM-based maps
  mislocate pilot jamming.
\item We introduce a DMRS residual for single-antenna receivers, the per-PRB energy of the
  pilot estimates outside the delay subspace, and derive its exact distribution with and
  without Gaussian jamming. The threshold for a given false-alarm rate follows from the
  allocation alone, and the false-alarm rate stays at its design value at every SNR.
\item We fuse power, residual and EVM in a small U-Net and evaluate it at 1\% per-cell false
  alarm against the threshold detectors, with matched ablations, a scheduling-aware attacker,
  unseen channel conditions, and network size and CPU latency.
\item We show how the map should be used by the receiver: as a statement of \emph{where} to
  estimate the jamming power, which then weights pilots and log-likelihood ratios (LLRs), rather
  than as a hard erasure mask~\cite{li2012robust}.
\end{itemize}

% ---------------------------------------------------------------------------------------------
\section{System and Threat Model}

We consider one PUSCH slot of 14 OFDM symbols on a carrier of 24 PRBs with 30\,kHz
subcarrier spacing~\cite{gpp20265g}. The UE has a single antenna, one layer, QPSK with LDPC
coding, and a contiguous allocation $\mathcal{A}$ of $n \ge 6$ PRBs. One front-loaded type-1
DMRS symbol carries pilots on every second subcarrier, so the allocation has $N = 6n$ pilots
of constant power $P_p$. The gNB has one antenna. The signal on subcarrier $k$ of symbol $l$
is
\begin{equation}
  y_{k,l} = h_{k,l}\, x_{k,l} + g_{k,l}\, z_{k,l} + w_{k,l},
  \label{eq:model}
\end{equation}
where $h$ and $g$ are the frequency responses of the UE's and the jammer's 3GPP TDL
channels~\cite{3gpp38901}, $x$ the UE's symbols (unit energy on data), $w \sim
\mathcal{CN}(0, N_0)$ the noise, and $z$ the jamming signal.

\emph{Threat model.} The jammer knows the numerology, the slot timing and the DMRS position,
but neither the data, the channels, the allocation nor the detector. It transmits
$z_{k,l} \sim \mathcal{CN}(0, S_J)$ on a set $\mathcal{C}$ of \emph{cells} (one PRB by one OFDM
symbol): a block of PRBs on the DMRS symbol only (\emph{pilot}), on the 13 other symbols
(\emph{data}), or on all symbols (\emph{partial-band}). Its strength is the jammer-to-signal
ratio per jammed resource element, $\mathrm{JSR} = S_J$; the jammer-to-noise ratio is
$\rho = S_J / N_0$. A scheduling-aware variant jams exactly the scheduled PRBs. The detector
outputs $\hat{Y} \in \{0,1\}^{24 \times 14}$ per slot; the label is $Y_{b,l} = 1$ if and only if
$(b,l) \in \mathcal{C}$, i.e., the attacked cells, not the cells whose symbols are degraded.

% ---------------------------------------------------------------------------------------------
\section{Receiver Features and the DMRS Residual}

All features are computed per cell from the receiver's own quantities and divided by $N_0$:
\begin{itemize}
\item \emph{Power} $P_{b,l} = \frac{1}{12 N_0} \sum_{k \in b} |y_{k,l}|^2$, on the whole carrier.
\item \emph{EVM} $E_{b,l} = \sum_{k\in b} |\hat{x}_{k,l} - \mathcal{Q}(\hat{x}_{k,l})|^2 /
  \sum_{k\in b} \nu_{k,l}$ on scheduled data cells, with $\hat{x}$ the LMMSE-equalized symbols,
  $\mathcal{Q}$ the nearest QPSK point and $\nu$ the equalizer's effective noise variance.
\item \emph{DMRS residual} $R_b$ on the DMRS cell of every scheduled PRB, defined next.
\end{itemize}
Power sees every attack but is masked by the UE's own fading. EVM sees jammed data, but it
cannot be computed on the DMRS symbol and, under pilot jamming, it rises on the unjammed data
symbols of the attacked PRBs: in our simulation (TDL-A, SNR 20\,dB, JSR 0\,dB on the pilots of
4 of 12 PRBs) the mean squared error of the equalized data symbols of those PRBs grows from 0.01
to 0.72, although the jammer never transmitted on them. EVM thus reports the right PRBs at the
wrong symbols. The residual addresses this blind spot.

\subsection{The residual}
Let $\tilde{h}_p = y_p / x_p$ be the least-squares estimates at the $N$ pilots. The channel
responses with delays in $[0, \tau_{\max}]$ lie, up to a fraction $10^{-6}$ of their energy, in
the span of the $L$ dominant eigenvectors $\mathbf{U} \in \mathbb{C}^{N \times L}$ of the
correlation matrix $[\mathbf{C}]_{pq} = e^{-j\pi\Delta f_{pq}\tau_{\max}}
\operatorname{sinc}(\Delta f_{pq}\tau_{\max})$ of a uniform delay profile ($\tau_{\max} =
3\,\mu$s, $L = 18$ for $n = 12$). With $\Pperp = \mathbf{I} - \mathbf{U}\mathbf{U}^H$, the
residual $\mathbf{r} = \Pperp \tilde{\mathbf{h}}$ removes the UE's channel and keeps noise and
jamming. For PRB $b$ with pilot set $\mathcal{B}$,
\begin{equation}
  R_b = \frac{\|\mathbf{r}_{\mathcal{B}}\|^2}{\sigma_p^2 \tr\big([\Pperp]_{\mathcal{B}\mathcal{B}}\big)},
  \qquad \sigma_p^2 = N_0 / P_p.
  \label{eq:residual}
\end{equation}

\begin{proposition}\label{prop:null}
Without jamming, and if the UE's channel at the pilots lies in the span of $\mathbf{U}$,
\begin{equation}
  R_b = \sum_{i=1}^{6} \lambda_{b,i}\, \epsilon_i, \quad \epsilon_i \overset{\text{i.i.d.}}{\sim}
  \mathrm{Exp}(1), \quad \lambda_{b,i} = \frac{\mu_i([\Pperp]_{\mathcal{B}\mathcal{B}})}
  {\tr([\Pperp]_{\mathcal{B}\mathcal{B}})},
  \label{eq:null}
\end{equation}
where $\mu_i(\cdot)$ are eigenvalues. The distribution depends on neither the channel, $N_0$
nor the pilot sequence, only on $n$ and the position of $b$ in the allocation.
\end{proposition}
\begin{IEEEproof}
$\Pperp \mathbf{h} = \mathbf{0}$, so $\mathbf{r}_{\mathcal{B}} = [\Pperp \tilde{\mathbf{w}}]_{\mathcal{B}}$
with $\tilde{w}_p = w_p / x_p \sim \mathcal{CN}(0, \sigma_p^2)$ i.i.d. ($|x_p|^2 = P_p$). Since
$\Pperp$ is a Hermitian projector, $\mathbf{r}_{\mathcal{B}} \sim \mathcal{CN}(\mathbf{0},
\sigma_p^2 [\Pperp]_{\mathcal{B}\mathcal{B}})$. Writing it in the eigenbasis of
$[\Pperp]_{\mathcal{B}\mathcal{B}}$ gives independent components whose squared magnitudes are
exponential with means $\sigma_p^2 \mu_i$.
\end{IEEEproof}

The threshold $\eta_b(\alpha)$ with $\Pr\{R_b > \eta_b\} = \alpha$ therefore follows from the
allocation alone, and the test $R_b > \eta_b(\alpha)$ is a constant-false-alarm-rate (CFAR)
test with respect to SNR and channel. In practice the condition $\Pperp\mathbf{h} = \mathbf{0}$
holds approximately: a single path at $\tau_{\max}$ leaves about $1.4\cdot 10^{-4}$ of its
energy in the residual of a band-edge PRB (at 2.6\,$\mu$s, the last tap of TDL-C, $4\cdot
10^{-6}$), which is negligible against the noise up to about 25\,dB SNR; beyond that, $\tau_{\max}$ or the retained energy fraction must grow. The tail of~\eqref{eq:null} is that of a chain of
exponential stages with rates $1/\lambda_{b,i}$, computed exactly as
$\Pr\{R_b > t\} = \mathbf{e}_1^T \exp(\mathbf{Q}_b t) \mathbf{1}$, with $\mathbf{Q}_b$ the
bidiagonal generator ($-1/\lambda_{b,i}$ on the diagonal, $1/\lambda_{b,i}$ above it); unlike
the partial-fraction form this is accurate for (nearly) equal weights. For $n = 12$ the 1\%
thresholds range from 2.31 in the band center to 2.56 at the band edges, where the subspace
vectors concentrate and fewer degrees of freedom remain.

\begin{proposition}\label{prop:alt}
Under Gaussian jamming of the pilots in a set $\mathcal{J}$ with jammer-to-noise ratio $\rho$
and jammer channel gains $|g_p|^2$, \eqref{eq:null} holds with weights
$\mu_i\big([\Pperp(\mathbf{I} + \rho\mathbf{G})\Pperp]_{\mathcal{B}\mathcal{B}}\big) /
\tr([\Pperp]_{\mathcal{B}\mathcal{B}})$, where $\mathbf{G} = \diag(|g_p|^2 \mathbb{I}[p \in
\mathcal{J}])$.
\end{proposition}
The proof is the same with $\tilde{\mathbf{w}}$ replaced by the noise plus $g_p z_p / x_p$.
Proposition~\ref{prop:alt} gives the detection probability of a jammed PRB and, for $b \notin
\mathcal{J}$, the false alarms caused by leakage: the projection spreads part of the jamming
energy onto neighbouring pilots. Averaging over the jammer's fading gives $P_d(\rho)$.

% ---------------------------------------------------------------------------------------------
\section{Fusion Detector}

The three features (in dB, scaled by 1/10) and three masks (allocation, DMRS cells, data cells)
form a $6 \times 24 \times 14$ input. A two-level U-Net~\cite{ronneberger2015unet} with $c$,
$2c$ and $4c$ channels ($c = 12$: 66\,553 parameters) outputs one logit per cell. It is trained
with a class-weighted binary cross-entropy for 15 epochs on 40\,000 slots, keeping the epoch
with the lowest validation loss. Every detector, learned or not, is thresholded so that it
flags 1\% of the non-jammed validation cells it scores. Ablations use the same network on
subsets of the inputs (power and EVM; power only; EVM only). The threshold baselines are the
EVM per symbol and resource block~\cite{ornek2024advanced}, the residual (the single-antenna
counterpart of the projected energy detector of~\cite{mehrabian2026robust}), the residual with
the analytic threshold of Proposition~\ref{prop:null}, and a per-cell power
test~\cite{li2012robust}.

% ---------------------------------------------------------------------------------------------
\section{Using the Map in the Receiver}
\label{sec:receiver}

The map can protect decoding in two ways. The classic \emph{hard} use~\cite{li2012robust}
drops the pilots of flagged PRBs, interpolates the channel from the others, and erases the LLRs
of flagged data cells. Every false alarm then removes clean information. We propose a
\emph{soft} use in which the map says where to estimate the jamming power. The channel is
interpolated by an LMMSE filter with the same delay model as the residual,
\begin{equation}
  \hat{h}_k = \mathbf{c}_k^H (\mathbf{C} + \mathbf{D})^{-1} \tilde{\mathbf{h}}, \quad
  e_k = 1 - \mathbf{c}_k^H (\mathbf{C} + \mathbf{D})^{-1} \mathbf{c}_k,
\end{equation}
where $\mathbf{c}_k$ is the correlation of subcarrier $k$ with the pilots and
$\mathbf{D} = \sigma_p^2 \diag(\delta_p)$, with $\delta_p = \max(R_{b(p)}, 1)$ if the DMRS cell
of the pilot's PRB is flagged and $\delta_p = 1$ otherwise; for a jammed PRB, $\mathbb{E}[R_b]$ is roughly $1 + \rho|g|^2$. For the data, the jamming power on PRB $b$ is estimated from the excess
received energy of its flagged data resource elements $\mathcal{F}_b$,
\begin{equation}
  \hat{S}_b = \max\Big(0, \frac{1}{|\mathcal{F}_b|} \sum_{(k,l) \in \mathcal{F}_b}
  |y_{k,l}|^2 - |\hat{h}_k|^2 - e_k - N_0 \Big),
\end{equation}
and added to the noise variance of those symbols: for one antenna, the unbiased LMMSE output
has effective noise $(e_k + N_0 + \hat{S}_b)/|\hat{h}_k|^2$, which scales the LLRs. Strong
jamming drives them to zero, as erasure would; after a false alarm $\hat{S}_b \approx 0$ and
$R_b \approx 1$, and nothing is lost.

A receiver without a map can only use its DMRS: we compare with a \emph{DMRS-only} receiver
that flags PRBs whose residual exceeds $\eta_b(1\%)$, weights their pilots as above, and,
assuming the interference is stationary over the slot, adds $N_0(R_b - 1)$ to the noise of all
their data symbols. It cannot see data jamming and penalizes the clean data of pilot-jammed
PRBs.

% ---------------------------------------------------------------------------------------------
\section{Results}
\label{sec:results}

\subsection{Setup}
The link is simulated with Sionna~\cite{hoydis2022sionna,nvidia2026sionna} (PyTorch backend):
24-PRB carrier at 3.5\,GHz, PUSCH with MCS~5 (QPSK, rate $\approx 0.37$), random contiguous
allocations of 6--24 PRBs, TDL-A (100\,ns) and TDL-C (300\,ns) channels at up to 3\,m/s, SNR
(per data resource element) drawn from 5--25\,dB. The jammer hits a block of 2--12 PRBs anywhere
on the carrier with JSR $\in [-15, 10]$\,dB. Training, validation and test sets (40\,000, 5\,000
and 31\,000 slots) use disjoint seeds; the test set is a grid over channel, attack and JSR
($-15$ to $+9$\,dB in 3\,dB steps) plus 4\,000 benign slots. Each U-Net is trained with three
seeds; we report mean $\pm$ standard deviation. Detection is measured per cell: $P_d$ over the
jammed cells inside the allocation, and the cell-level F1 score, at thresholds that flag 1\% of
the non-jammed validation cells a detector scores.

\subsection{The residual follows its theory}
Fig.~\ref{fig:theory}(a) applies the analytic 1\% threshold of Proposition~\ref{prop:null},
computed without any data, to every scheduled DMRS cell of the benign test slots. The measured
false-alarm rate is 0.91\% (TDL-A) and 0.92\% (TDL-C) and does not move between 5 and 25\,dB
SNR. Fig.~\ref{fig:theory}(b) compares Proposition~\ref{prop:alt} with the simulation for pilot
jamming of 4 of 12 PRBs: inside the block, the formula with block-Rayleigh jammer fading
predicts $P_d = 0.29, 0.69, 0.93$ at JNR 0, 6, 12\,dB, and the simulation gives 0.30,
0.71--0.74, 0.93--0.94; for the unjammed PRB next to the block it predicts the leakage-induced
false alarms (0.16 at 12\,dB, 0.50 at 18\,dB; simulated 0.17 and 0.54). The assumption of
Proposition~\ref{prop:null} matters: on TDL-C with 1000\,ns delay spread, whose delays reach
8.7\,$\mu$s $> \tau_{\max}$, part of the UE's channel stays in the residual and the false-alarm
rate climbs to 61\%.

\begin{figure*}[t]
\centering
\includegraphics[width=0.92\textwidth]{theory/residual_theory.pdf}
\caption{DMRS residual against Propositions~\ref{prop:null} and~\ref{prop:alt}. (a) False
alarms at the analytic 1\% threshold in benign slots. (b) Probability that the residual of a
PRB exceeds that threshold under pilot jamming of PRBs 4--7 of a 12-PRB allocation (SNR 10\,dB):
lines are the formula, markers the link simulation.}
\label{fig:theory}
\end{figure*}

\subsection{Detection}
Table~\ref{tab:detection} and Fig.~\ref{fig:pd} summarize detection. The fused U-Net reaches an
overall F1 of $0.951 \pm 0.002$. The residual is what makes weak pilot jamming detectable: at
$-15$\,dB JSR it raises $P_d$ from 0.43 (power and EVM) to 0.62, and the pilot-attack F1 from
0.757 to 0.834; for data and partial-band jamming the two networks are equal, as they should be,
since these attacks leave the DMRS alone. Every threshold detector sees only part of the
problem. The EVM threshold~\cite{ornek2024advanced} cannot score the DMRS symbol at all and,
under pilot jamming, flags the unjammed data symbols of the attacked PRBs instead (F1 0 on
pilot attacks). The residual tests are blind to data jamming. The per-cell power
test~\cite{li2012robust}, even relative to the slot's median scheduled power, cannot separate
jamming from the UE's own frequency-selective fading (F1 0.12). The power-only U-Net, given the
same power map, reaches 0.913: it compares each cell with the other symbols of the same PRB,
which share its channel at low mobility.

\begin{table}[t]
\caption{Cell-level F1 at 1\% per-cell false alarm and $P_d$ at JSR $-15$\,dB (scheduled cells;
U-Nets: mean $\pm$ std over three seeds)}
\label{tab:detection}
\centering
\setlength{\tabcolsep}{3pt}
\begin{tabular}{@{}lcccccc@{}}
\toprule
 & \multicolumn{4}{c}{F1} & \multicolumn{2}{c}{$P_d$, $-15$\,dB} \\
\cmidrule(lr){2-5}\cmidrule(l){6-7}
Detector & pilot & data & partial & all & pilot & data \\
\midrule
U-Net P+R+E & \textbf{0.834} & \textbf{0.957} & \textbf{0.961} & \textbf{0.951} & \textbf{0.62} & \textbf{0.78} \\
U-Net P+E & 0.757 & 0.952 & 0.960 & 0.946 & 0.43 & 0.74 \\
U-Net P & 0.817 & 0.920 & 0.918 & 0.913 & 0.25 & 0.55 \\
U-Net E & 0.459 & 0.792 & 0.782 & 0.772 & 0.14 & 0.47 \\
EVM threshold~\cite{ornek2024advanced} & 0 & 0.324 & 0.319 & 0.303 & 0 & 0 \\
Residual threshold & 0.584 & 0 & 0.057 & 0.058 & 0.10 & 0 \\
Residual, exact $p$-value & 0.588 & 0 & 0.058 & 0.059 & 0.10 & 0 \\
Power threshold~\cite{li2012robust} & 0.209 & 0.128 & 0.106 & 0.120 & 0.01 & 0.02 \\
\bottomrule
\end{tabular}
\end{table}

\begin{figure*}[t]
\centering
\includegraphics[width=0.92\textwidth]{eval/test_pd_vs_jsr_scheduled.pdf}
\caption{Detection probability at 1\% per-cell false alarm against JSR, jammed cells inside the
allocation, TDL-A and TDL-C (bands: $\pm 1$ std over training seeds).}
\label{fig:pd}
\end{figure*}

\subsection{Robustness}
\emph{Scheduling-aware jammer.} When the jammer hits exactly the scheduled PRBs, an attacker no
network saw in training, there is no unjammed neighbouring PRB to compare with. For pilot
jamming the fused detector is hardly affected ($P_d = 0.59$ at $-15$\,dB against 0.62 in the
test set, F1 $0.880 \pm 0.011$), because the residual is an absolute measure; without it $P_d$
falls from 0.43 to 0.35. For data jamming every network loses some sensitivity (overall F1 0.934
for P+R+E, 0.937 for P+E): part of what they learned is contrast between jammed and unjammed
PRBs.

\emph{Unseen channels and waveforms.} Table~\ref{tab:unseen} keeps the thresholds from the
matched validation set and changes one condition at a time. A constant-envelope QPSK jammer
changes nothing for the fused detector. Delays beyond $\tau_{\max}$ violate the assumption of
Proposition~\ref{prop:null} and raise its false alarms; at 30\,m/s the single DMRS symbol keeps
the residual intact (1.2\% false alarms), but the channel changes within the slot and the EVM of
the late symbols grows without jamming. The power-only network is robust to both. Training on more diverse channels removes most of the loss
(last rows of Table~\ref{tab:unseen}): with batches cycling through the nominal channels, long
delay spreads and speeds up to 30\,m/s, the fused detector reaches 0.921 at 1000\,ns and 0.946 at
30\,m/s with false-alarm rates of 1.7\% and 1.0\%, and 0.955 in matched conditions. These channel
types were part of the mix (with independent realizations), so this shows that diversity is
sufficient, not that arbitrary channels are handled.

\begin{table}[t]
\caption{F1 / false-alarm rate on unseen conditions (thresholds from the matched validation
set; mean over three seeds)}
\label{tab:unseen}
\centering
\setlength{\tabcolsep}{3pt}
\begin{tabular}{@{}lccc@{}}
\toprule
Condition & U-Net P+R+E & U-Net P+E & U-Net P \\
\midrule
Matched test & \textbf{0.951} / 1.2\% & 0.946 / 1.3\% & 0.913 / 1.2\% \\
QPSK jammer & \textbf{0.951} / 1.1\% & 0.942 / 1.3\% & 0.905 / 1.2\% \\
TDL-C, 1000\,ns & 0.818 / 6.1\% & 0.701 / 15.7\% & \textbf{0.923} / 1.3\% \\
30\,m/s & 0.865 / 5.1\% & 0.835 / 6.7\% & \textbf{0.912} / 1.5\% \\
\midrule
\multicolumn{4}{@{}l}{\emph{Trained on mixed channels:}} \\
Matched test & \textbf{0.955} / 0.7\% & 0.948 / 0.7\% & 0.914 / 1.0\% \\
TDL-C, 1000\,ns & 0.921 / 1.7\% & 0.911 / 2.1\% & \textbf{0.928} / 1.0\% \\
30\,m/s & \textbf{0.946} / 1.0\% & 0.939 / 1.0\% & 0.917 / 1.1\% \\
\bottomrule
\end{tabular}
\end{table}

\subsection{Size and latency}
A 30\,kHz slot lasts 0.5\,ms. On one core of an Apple M1 Pro CPU, feature extraction takes
0.26\,ms per slot and the 66k-parameter U-Net 0.49\,ms, more than a slot in total. The network
is larger than it needs to be: with the first level narrowed from 12 to 4 channels (7\,657
parameters) the overall F1 drops only from 0.951 to $0.943 \pm 0.002$, mostly through weak pilot
jamming ($P_d$ 0.55 instead of 0.62 at $-15$\,dB), and the complete detector runs in 0.53\,ms in
eager PyTorch and 0.43\,ms as a frozen graph with batch normalization folded into the
convolutions. With 16\,933 parameters (F1 0.948) it still fits, at 0.48\,ms.

\subsection{Using the map in the receiver}
Fig.~\ref{fig:recovery} shows the BLER of the receivers of Section~\ref{sec:receiver} on TDL-C at 15\,dB SNR,
with the UE on 12 PRBs and the jammer on 4 of them. All receivers decode the same slots; each
runs until 100 block errors or 3\,200 slots. Without protection, jamming at $+9$\,dB JSR breaks
the link (BLER 0.33 for pilot, 0.84 for data, 0.90 for partial-band jamming). The DMRS-only
receiver recovers pilot and partial-band jamming (0.008 and 0.011) but cannot see data jamming
(0.84): the map's resolution in time is what makes data jamming recoverable. The power-test map
used hard~\cite{li2012robust} helps only against strong jamming (0.18 for data jamming at
$+9$\,dB). With the fused map used softly, the BLER stays at or below 0.001 for every attack and
JSR, and at 0 without a jammer, the same as with the true map.

The difference between hard and soft use appears at weak jamming, where the unprotected link
decodes. Used hard, the fused map raises the BLER to 0.12 at $-15$\,dB data or partial-band
jamming, because false alarms drop clean pilots and erase clean data; even the true map costs
about 0.002, since erasing a third of the allocation discards symbols that weak jamming barely
damaged. Used softly, neither costs anything. With soft use, the map without the residual is as
good as the fused one: weak jamming that it misses also barely hurts decoding, and the receiver
weights the pilots by their residual anyway. When every pilot of the allocation is jammed, no
clean pilot is left to interpolate from; weighting all of them by their residual still lets the
LMMSE filter average across the band, which lowers the BLER at $+3$\,dB from 0.41 to 0.04.
Results on TDL-A are similar.

\begin{figure*}[t]
\centering
\includegraphics[width=0.92\textwidth]{recovery/recovery_C.pdf}
\caption{BLER against JSR on TDL-C at 15\,dB SNR (UE on 12 PRBs, jammer on 4 of them), all
receivers decoding the same slots. Hollow triangles: no block error, drawn at the 95\% upper
bound; bands: 95\% confidence intervals.}
\label{fig:recovery}
\end{figure*}

% ---------------------------------------------------------------------------------------------
\section{Conclusion}

A gNB can localize resource-selective jamming at PRB-by-symbol resolution from what its
receiver already computes. The DMRS residual turns the known structure of physical channels into
a test whose false-alarm rate is fixed by design: its distribution without jamming depends only on
the allocation, which gives an analytic threshold and closed-form detection and leakage
probabilities, confirmed by the link simulation. Fused with power and EVM in a U-Net of a few
thousand parameters, it localizes pilot, data and partial-band jamming at 1\% per-cell false
alarm, including weak pilot jamming that EVM-based maps place on the wrong symbols, within one
slot of CPU time. Used to decide where the receiver estimates the jamming power, rather than what
it erases, the map protects the link against strong jamming at no measurable cost under weak jamming
or false alarms.

The study is limited to one receive antenna, a Gaussian jammer aligned with the slot grid,
perfect synchronization and known noise power. The residual generalizes naturally to multiple
antennas (a space--frequency subspace) and to additional DMRS symbols, which would also address
fast users; delay spreads beyond $\tau_{\max}$ require a wider subspace or training on such
channels, as our mixed-training results show.

\bibliographystyle{IEEEtran}
\bibliography{../docs/references}

\end{document}
